% Options for packages loaded elsewhere
\PassOptionsToPackage{unicode}{hyperref}
\PassOptionsToPackage{hyphens}{url}
%
\documentclass[
]{article}
\usepackage{amsmath,amssymb}
\usepackage{iftex}
\ifPDFTeX
  \usepackage[T1]{fontenc}
  \usepackage[utf8]{inputenc}
  \usepackage{textcomp} % provide euro and other symbols
\else % if luatex or xetex
  \usepackage{unicode-math} % this also loads fontspec
  \defaultfontfeatures{Scale=MatchLowercase}
  \defaultfontfeatures[\rmfamily]{Ligatures=TeX,Scale=1}
\fi
\usepackage{lmodern}
\ifPDFTeX\else
  % xetex/luatex font selection
\fi
% Use upquote if available, for straight quotes in verbatim environments
\IfFileExists{upquote.sty}{\usepackage{upquote}}{}
\IfFileExists{microtype.sty}{% use microtype if available
  \usepackage[]{microtype}
  \UseMicrotypeSet[protrusion]{basicmath} % disable protrusion for tt fonts
}{}
\makeatletter
\@ifundefined{KOMAClassName}{% if non-KOMA class
  \IfFileExists{parskip.sty}{%
    \usepackage{parskip}
  }{% else
    \setlength{\parindent}{0pt}
    \setlength{\parskip}{6pt plus 2pt minus 1pt}}
}{% if KOMA class
  \KOMAoptions{parskip=half}}
\makeatother
\usepackage{xcolor}
\usepackage[margin=1in]{geometry}
\usepackage{color}
\usepackage{fancyvrb}
\newcommand{\VerbBar}{|}
\newcommand{\VERB}{\Verb[commandchars=\\\{\}]}
\DefineVerbatimEnvironment{Highlighting}{Verbatim}{commandchars=\\\{\}}
% Add ',fontsize=\small' for more characters per line
\usepackage{framed}
\definecolor{shadecolor}{RGB}{248,248,248}
\newenvironment{Shaded}{\begin{snugshade}}{\end{snugshade}}
\newcommand{\AlertTok}[1]{\textcolor[rgb]{0.94,0.16,0.16}{#1}}
\newcommand{\AnnotationTok}[1]{\textcolor[rgb]{0.56,0.35,0.01}{\textbf{\textit{#1}}}}
\newcommand{\AttributeTok}[1]{\textcolor[rgb]{0.13,0.29,0.53}{#1}}
\newcommand{\BaseNTok}[1]{\textcolor[rgb]{0.00,0.00,0.81}{#1}}
\newcommand{\BuiltInTok}[1]{#1}
\newcommand{\CharTok}[1]{\textcolor[rgb]{0.31,0.60,0.02}{#1}}
\newcommand{\CommentTok}[1]{\textcolor[rgb]{0.56,0.35,0.01}{\textit{#1}}}
\newcommand{\CommentVarTok}[1]{\textcolor[rgb]{0.56,0.35,0.01}{\textbf{\textit{#1}}}}
\newcommand{\ConstantTok}[1]{\textcolor[rgb]{0.56,0.35,0.01}{#1}}
\newcommand{\ControlFlowTok}[1]{\textcolor[rgb]{0.13,0.29,0.53}{\textbf{#1}}}
\newcommand{\DataTypeTok}[1]{\textcolor[rgb]{0.13,0.29,0.53}{#1}}
\newcommand{\DecValTok}[1]{\textcolor[rgb]{0.00,0.00,0.81}{#1}}
\newcommand{\DocumentationTok}[1]{\textcolor[rgb]{0.56,0.35,0.01}{\textbf{\textit{#1}}}}
\newcommand{\ErrorTok}[1]{\textcolor[rgb]{0.64,0.00,0.00}{\textbf{#1}}}
\newcommand{\ExtensionTok}[1]{#1}
\newcommand{\FloatTok}[1]{\textcolor[rgb]{0.00,0.00,0.81}{#1}}
\newcommand{\FunctionTok}[1]{\textcolor[rgb]{0.13,0.29,0.53}{\textbf{#1}}}
\newcommand{\ImportTok}[1]{#1}
\newcommand{\InformationTok}[1]{\textcolor[rgb]{0.56,0.35,0.01}{\textbf{\textit{#1}}}}
\newcommand{\KeywordTok}[1]{\textcolor[rgb]{0.13,0.29,0.53}{\textbf{#1}}}
\newcommand{\NormalTok}[1]{#1}
\newcommand{\OperatorTok}[1]{\textcolor[rgb]{0.81,0.36,0.00}{\textbf{#1}}}
\newcommand{\OtherTok}[1]{\textcolor[rgb]{0.56,0.35,0.01}{#1}}
\newcommand{\PreprocessorTok}[1]{\textcolor[rgb]{0.56,0.35,0.01}{\textit{#1}}}
\newcommand{\RegionMarkerTok}[1]{#1}
\newcommand{\SpecialCharTok}[1]{\textcolor[rgb]{0.81,0.36,0.00}{\textbf{#1}}}
\newcommand{\SpecialStringTok}[1]{\textcolor[rgb]{0.31,0.60,0.02}{#1}}
\newcommand{\StringTok}[1]{\textcolor[rgb]{0.31,0.60,0.02}{#1}}
\newcommand{\VariableTok}[1]{\textcolor[rgb]{0.00,0.00,0.00}{#1}}
\newcommand{\VerbatimStringTok}[1]{\textcolor[rgb]{0.31,0.60,0.02}{#1}}
\newcommand{\WarningTok}[1]{\textcolor[rgb]{0.56,0.35,0.01}{\textbf{\textit{#1}}}}
\usepackage{longtable,booktabs,array}
\usepackage{calc} % for calculating minipage widths
% Correct order of tables after \paragraph or \subparagraph
\usepackage{etoolbox}
\makeatletter
\patchcmd\longtable{\par}{\if@noskipsec\mbox{}\fi\par}{}{}
\makeatother
% Allow footnotes in longtable head/foot
\IfFileExists{footnotehyper.sty}{\usepackage{footnotehyper}}{\usepackage{footnote}}
\makesavenoteenv{longtable}
\usepackage{graphicx}
\makeatletter
\newsavebox\pandoc@box
\newcommand*\pandocbounded[1]{% scales image to fit in text height/width
  \sbox\pandoc@box{#1}%
  \Gscale@div\@tempa{\textheight}{\dimexpr\ht\pandoc@box+\dp\pandoc@box\relax}%
  \Gscale@div\@tempb{\linewidth}{\wd\pandoc@box}%
  \ifdim\@tempb\p@<\@tempa\p@\let\@tempa\@tempb\fi% select the smaller of both
  \ifdim\@tempa\p@<\p@\scalebox{\@tempa}{\usebox\pandoc@box}%
  \else\usebox{\pandoc@box}%
  \fi%
}
% Set default figure placement to htbp
\def\fps@figure{htbp}
\makeatother
\setlength{\emergencystretch}{3em} % prevent overfull lines
\providecommand{\tightlist}{%
  \setlength{\itemsep}{0pt}\setlength{\parskip}{0pt}}
\setcounter{secnumdepth}{5}
% definitions for citeproc citations
\NewDocumentCommand\citeproctext{}{}
\NewDocumentCommand\citeproc{mm}{%
  \begingroup\def\citeproctext{#2}\cite{#1}\endgroup}
\makeatletter
 % allow citations to break across lines
 \let\@cite@ofmt\@firstofone
 % avoid brackets around text for \cite:
 \def\@biblabel#1{}
 \def\@cite#1#2{{#1\if@tempswa , #2\fi}}
\makeatother
\newlength{\cslhangindent}
\setlength{\cslhangindent}{1.5em}
\newlength{\csllabelwidth}
\setlength{\csllabelwidth}{3em}
\newenvironment{CSLReferences}[2] % #1 hanging-indent, #2 entry-spacing
 {\begin{list}{}{%
  \setlength{\itemindent}{0pt}
  \setlength{\leftmargin}{0pt}
  \setlength{\parsep}{0pt}
  % turn on hanging indent if param 1 is 1
  \ifodd #1
   \setlength{\leftmargin}{\cslhangindent}
   \setlength{\itemindent}{-1\cslhangindent}
  \fi
  % set entry spacing
  \setlength{\itemsep}{#2\baselineskip}}}
 {\end{list}}
\usepackage{calc}
\newcommand{\CSLBlock}[1]{\hfill\break\parbox[t]{\linewidth}{\strut\ignorespaces#1\strut}}
\newcommand{\CSLLeftMargin}[1]{\parbox[t]{\csllabelwidth}{\strut#1\strut}}
\newcommand{\CSLRightInline}[1]{\parbox[t]{\linewidth - \csllabelwidth}{\strut#1\strut}}
\newcommand{\CSLIndent}[1]{\hspace{\cslhangindent}#1}
\usepackage{bookmark}
\IfFileExists{xurl.sty}{\usepackage{xurl}}{} % add URL line breaks if available
\urlstyle{same}
\hypersetup{
  pdftitle={Dynamic Reporting with RMarkdown},
  hidelinks,
  pdfcreator={LaTeX via pandoc}}

\title{Dynamic Reporting with RMarkdown}
\usepackage{etoolbox}
\makeatletter
\providecommand{\subtitle}[1]{% add subtitle to \maketitle
  \apptocmd{\@title}{\par {\large #1 \par}}{}{}
}
\makeatother
\subtitle{Universal Template (PDF, HTML, Word)}
\author{}
\date{\vspace{-2.5em}2025-12-30}

\begin{document}
\maketitle
\begin{abstract}
This document serves as a universal template. It is designed to be
rendered into HTML, PDF, or MS Word without requiring changes to the
code. It demonstrates automatic table formatting and bibliography
handling.
\end{abstract}

{
\setcounter{tocdepth}{2}
\tableofcontents
}
\newpage

\section{Introduction}\label{introduction}

This report demonstrates how to write \textbf{one} RMarkdown file that
compiles cleanly to PDF, HTML, and Word. We avoid using
\texttt{kableExtra} or raw \texttt{\textless{}div\textgreater{}} tags,
relying instead on pure Pandoc Markdown.

\section{Tables}\label{tables}

To ensure tables look good in Word, PDF, and HTML simultaneously, we use
\texttt{knitr::kable} with the markdown format.

\begin{Shaded}
\begin{Highlighting}[]
\FunctionTok{head}\NormalTok{(mtcars[, }\DecValTok{1}\SpecialCharTok{:}\DecValTok{4}\NormalTok{]) }\SpecialCharTok{|\textgreater{}} 
  \FunctionTok{kable}\NormalTok{(}\AttributeTok{caption =} \StringTok{"A Universal Table Example"}\NormalTok{, }
        \AttributeTok{format =} \StringTok{"markdown"}\NormalTok{)  }\CommentTok{\# "markdown" format is safest for cross{-}compatibility}
\end{Highlighting}
\end{Shaded}

\begin{longtable}[]{@{}lrrrr@{}}
\caption{A Universal Table Example}\tabularnewline
\toprule\noalign{}
& mpg & cyl & disp & hp \\
\midrule\noalign{}
\endfirsthead
\toprule\noalign{}
& mpg & cyl & disp & hp \\
\midrule\noalign{}
\endhead
\bottomrule\noalign{}
\endlastfoot
Mazda RX4 & 21.0 & 6 & 160 & 110 \\
Mazda RX4 Wag & 21.0 & 6 & 160 & 110 \\
Datsun 710 & 22.8 & 4 & 108 & 93 \\
Hornet 4 Drive & 21.4 & 6 & 258 & 110 \\
Hornet Sportabout & 18.7 & 8 & 360 & 175 \\
Valiant & 18.1 & 6 & 225 & 105 \\
\end{longtable}

\section{Graphics}\label{graphics}

\subsection{Standard Vector Graphics}\label{standard-vector-graphics}

Standard graphics are rendered as vectors (in PDF) or high-quality
images (HTML/Word) by default.

Graphics can be embedded into the documents as shown below.

\subsubsection{Package: base}\label{package-base}

\begin{figure}

{\centering \includegraphics{Report_files/figure-latex/baseGrafik-1} 

}

\caption{This is a **base** graphic}\label{fig:baseGrafik}
\end{figure}

\subsubsection{Package: ggplot2}\label{package-ggplot2}

See the cheat sheet for ggplot2 graphics for details:
\url{https://posit.co/wp-content/uploads/2022/10/data-visualization-1.pdf}

\begin{figure}

{\centering \includegraphics{Report_files/figure-latex/ggplotGrafik-1} 

}

\caption{**ggplot** graphic. Captions with dynamic content are also possible, e.g., Date and Time: 2025-12-30 00:33:47.738028}\label{fig:ggplotGrafik}
\end{figure}

The use of chunk options to set captions allows for dynamic numbering
and list of figures generation. This works also for \texttt{base}
graphics (see above). Check the chunk option \texttt{fig.asp=0.5} that
is responsible for the altered aspect ratio.

\begin{Shaded}
\begin{Highlighting}[]
\FunctionTok{ggplot}\NormalTok{(iris, }\FunctionTok{aes}\NormalTok{(}\AttributeTok{x =}\NormalTok{ Sepal.Width, }\AttributeTok{y =}\NormalTok{ Petal.Width)) }\SpecialCharTok{+}
  \FunctionTok{geom\_point}\NormalTok{(}\FunctionTok{aes}\NormalTok{(}\AttributeTok{color =}\NormalTok{ Species)) }
\end{Highlighting}
\end{Shaded}

\begin{figure}

{\centering \includegraphics{Report_files/figure-latex/ggplotGrafikAspectChanged-1} 

}

\caption{**ggplot** graphic with different aspect ratio.}\label{fig:ggplotGrafikAspectChanged}
\end{figure}

\subsection{Large Datasets \&
Rasterization}\label{large-datasets-rasterization}

When plotting datasets with thousands of points (like the
\texttt{diamonds} dataset), vector graphics can make the final file huge
and slow to load.

To solve this, we use the chunk option
\texttt{dev=\textquotesingle{}png\textquotesingle{}} and
\texttt{dpi=300}. This forces this specific plot to render as a flat
image (raster), keeping the document size small while maintaining
transparency.

\begin{Shaded}
\begin{Highlighting}[]
\CommentTok{\# The \textquotesingle{}alpha\textquotesingle{} parameter adds transparency to reveal density}
\NormalTok{diamonds }\SpecialCharTok{\%\textgreater{}\%} 
\FunctionTok{ggplot}\NormalTok{(}\FunctionTok{aes}\NormalTok{(}\AttributeTok{x =}\NormalTok{ carat, }\AttributeTok{y =}\NormalTok{ price, }\AttributeTok{color =}\NormalTok{ cut)) }\SpecialCharTok{+}
  \FunctionTok{geom\_point}\NormalTok{(}\AttributeTok{alpha =} \FloatTok{0.1}\NormalTok{) }\SpecialCharTok{+} 
  \FunctionTok{labs}\NormalTok{(}\AttributeTok{title =} \StringTok{"Diamonds: Price vs Carat"}\NormalTok{,}
       \AttributeTok{subtitle =} \StringTok{"50,000+ points rendered as PNG to save space"}\NormalTok{,}
       \AttributeTok{x =} \StringTok{"Carat"}\NormalTok{, }
       \AttributeTok{y =} \StringTok{"Price"}\NormalTok{) }\SpecialCharTok{+}
  \FunctionTok{theme}\NormalTok{(}\AttributeTok{legend.position =} \StringTok{"bottom"}\NormalTok{)}
\end{Highlighting}
\end{Shaded}

\begin{figure}

{\centering \includegraphics{Report_files/figure-latex/heavyPlot-1} 

}

\caption{Figure 2: Rasterized plot (PNG) of 50,000+ points.}\label{fig:heavyPlot}
\end{figure}

\section{Summary Statistics}\label{summary-statistics}

We can also include raw text outputs, which are rendered as code blocks
in all formats.

\begin{Shaded}
\begin{Highlighting}[]
\FunctionTok{summary}\NormalTok{(mtcars}\SpecialCharTok{$}\NormalTok{mpg)}
\end{Highlighting}
\end{Shaded}

\begin{verbatim}
##    Min. 1st Qu.  Median    Mean 3rd Qu.    Max. 
##   10.40   15.43   19.20   20.09   22.80   33.90
\end{verbatim}

\section{Citations}\label{citations}

The information about the citations used are stored in \texttt{bibtex}
style in the additional file \texttt{references.bib} which the last line
of the YAML header is referencing to.

You can cite entries using the syntax \texttt{{[}@key{]}}.

We are using the \textbf{knitr} package (Xie 2025) for processing and
\textbf{ggplot2} (Wickham et al. 2025) for visualization. References
will be placed automatically at the end of the document or where we
specify the placement div.

There are alternative ways to cite:

\begin{itemize}
\tightlist
\item
  Indirect citation: This report was built using the knitr package (Xie
  2025).
\item
  Direct citation: Wickham et al. (2025) describes the grammar of
  graphics.
\item
  Multiple citations: RMarkdown is powerful (Allaire et al. 2025; Xie
  2025).
\end{itemize}

\section{Literature}\label{literature}

\phantomsection\label{refs}
\begin{CSLReferences}{1}{0}
\bibitem[\citeproctext]{ref-R-rmarkdown}
Allaire, JJ, Yihui Xie, Christophe Dervieux, Jonathan McPherson, Javier
Luraschi, Kevin Ushey, Aron Atkins, et al. 2025. \emph{Rmarkdown:
Dynamic Documents for r}. \url{https://github.com/rstudio/rmarkdown}.

\bibitem[\citeproctext]{ref-R-ggplot2}
Wickham, Hadley, Winston Chang, Lionel Henry, Thomas Lin Pedersen,
Kohske Takahashi, Claus Wilke, Kara Woo, Hiroaki Yutani, Dewey
Dunnington, and Teun van den Brand. 2025. \emph{Ggplot2: Create Elegant
Data Visualisations Using the Grammar of Graphics}.
\url{https://ggplot2.tidyverse.org}.

\bibitem[\citeproctext]{ref-R-knitr}
Xie, Yihui. 2025. \emph{Knitr: A General-Purpose Package for Dynamic
Report Generation in r}. \url{https://yihui.org/knitr/}.

\end{CSLReferences}

\section{Appendix}\label{appendix}

\subsection{Appendix A: System
Information}\label{appendix-a-system-information}

Because we placed the
\texttt{\textless{}div\ id="refs"\textgreater{}\textless{}/div\textgreater{}}
above, this Appendix will correctly appear \textbf{after} the references
in the final document.

\begin{Shaded}
\begin{Highlighting}[]
\FunctionTok{sessionInfo}\NormalTok{()}
\end{Highlighting}
\end{Shaded}

\begin{verbatim}
## R version 4.4.1 (2024-06-14 ucrt)
## Platform: x86_64-w64-mingw32/x64
## Running under: Windows 11 x64 (build 22621)
## 
## Matrix products: default
## 
## 
## locale:
## [1] LC_COLLATE=English_United States.utf8 
## [2] LC_CTYPE=English_United States.utf8   
## [3] LC_MONETARY=English_United States.utf8
## [4] LC_NUMERIC=C                          
## [5] LC_TIME=English_United States.utf8    
## 
## time zone: Europe/Paris
## tzcode source: internal
## 
## attached base packages:
## [1] stats     graphics  grDevices utils     datasets  methods   base     
## 
## other attached packages:
## [1] dplyr_1.1.4   ggplot2_3.5.2 knitr_1.50   
## 
## loaded via a namespace (and not attached):
##  [1] vctrs_0.6.5       cli_3.6.4         rlang_1.1.6       xfun_0.52        
##  [5] generics_0.1.3    labeling_0.4.3    glue_1.8.0        colorspace_2.1-1 
##  [9] htmltools_0.5.8.1 scales_1.3.0      rmarkdown_2.29    grid_4.4.1       
## [13] evaluate_1.0.3    munsell_0.5.1     tibble_3.2.1      fastmap_1.2.0    
## [17] yaml_2.3.10       lifecycle_1.0.4   compiler_4.4.1    pkgconfig_2.0.3  
## [21] rstudioapi_0.17.1 farver_2.1.2      digest_0.6.37     viridisLite_0.4.2
## [25] R6_2.6.1          tidyselect_1.2.1  pillar_1.10.2     magrittr_2.0.3   
## [29] withr_3.0.2       tools_4.4.1       gtable_0.3.6
\end{verbatim}

\end{document}
